
    %%%%%%%%%%%%%%%%%%%%%%%%%%%%%%%%%%%%%%%%%%%%%%%%%%%%%%%%%%%%%%%%%%%%%%%%%%%%%%%%
    %%% Classe do documento
    \documentclass[12pt, addpoints]{exam}
    \UseRawInputEncoding

    %%%%%%%%%%%%%%%%%%%%%%%%%%%%%%%%%%%%%%%%%%%%%%%%%%%%%%%%%%%%%%%%%%%%%%%%%%%%%%%%
    % preambulo.tex
    %
    % Arquivo de configuração de packages para uso o LaTeX, conforme minhas
    % preferências e modelos pessoais.
    %
    % Para maiores informações, visite:
    %    https://github.com/abrantesasf/latex
    %
    % NÃO ALTERE SE NÃO SOUBER O QUE ESTÁ FAZENDO!
    %%%%%%%%%%%%%%%%%%%%%%%%%%%%%%%%%%%%%%%%%%%%%%%%%%%%%%%%%%%%%%%%%%%%%%%%%%%%%%%%


    %%%%%%%%%%%%%%%%%%%%%%%%%%%%%%%%%%%%%%%%%%%%%%%%%%%%%%%%%%%%%%%%%%%%%%%%%%%%%%%%
    %%% Estruturas de controle:
    \input{static/utils/controles}

    %%%%%%%%%%%%%%%%%%%%%%%%%%%%%%%%%%%%%%%%%%%%%%%%%%%%%%%%%%%%%%%%%%%%%%%%%%%%%%%%
    %%% Compilação condicional em PDF:
    \input{static/utils/pdf}

    %%%%%%%%%%%%%%%%%%%%%%%%%%%%%%%%%%%%%%%%%%%%%%%%%%%%%%%%%%%%%%%%%%%%%%%%%%%%%%%%
    %%% Configurações de layout da página:
    \input{static/utils/layout}

    %%%%%%%%%%%%%%%%%%%%%%%%%%%%%%%%%%%%%%%%%%%%%%%%%%%%%%%%%%%%%%%%%%%%%%%%%%%%%%%%
    %%% Configurações para autores:
    \input{static/utils/autores}

    %%%%%%%%%%%%%%%%%%%%%%%%%%%%%%%%%%%%%%%%%%%%%%%%%%%%%%%%%%%%%%%%%%%%%%%%%%%%%%%%
    %%% Cabeçalho e rodapé:
    %%% Atenção! É MUITO DIFÍCIL ajustar apropriadamente os running headers
    %%% e running footers da primeira vez. Você pode confiar no LaTeX e deixar que
    %%% ele faça o ajuste automaticamente ou pode quebrar a cabeça e
    %%% tentar melhorar algo que o LaTeX já faz muito bem, mesmo que não fique
    %%% exatamente do jeito que você quer. O que você prefere? Se quiser ajustar
    %%% manualmente, descomente a linha abaixo e faça as configurações no arquivo
    %%% "utils/headings.tex".
    %\input{static/utils/headings}

    %%%%%%%%%%%%%%%%%%%%%%%%%%%%%%%%%%%%%%%%%%%%%%%%%%%%%%%%%%%%%%%%%%%%%%%%%%%%%%%%
    %%% Configuração para as fontes do documento:
    %%% Se você quiser usar fontes próprias ou do sistema, precisará ajustar as
    %%% configurações no arquivo "utils/fontes.tex".
    %%% ATENÇÃO: por padrão eu utilizo XeLaTeX com fontes proprietárias projetadas
    %%% por Matthew Butterick, adquiridas em https://mbtype.com, que não podem ser
    %%% distribuídas devido à licença de utilização. Se você usar XeLaTeX, você
    %%% DEVERÁ OBRIGATORIAMENTE ajustar as fontes no arquivo "utils/fontes.tex".
    \input{static/utils/fontes}

    %%%%%%%%%%%%%%%%%%%%%%%%%%%%%%%%%%%%%%%%%%%%%%%%%%%%%%%%%%%%%%%%%%%%%%%%%%%%%%%%
    %%% Sumário:
    \input{static/utils/toc}

    %%%%%%%%%%%%%%%%%%%%%%%%%%%%%%%%%%%%%%%%%%%%%%%%%%%%%%%%%%%%%%%%%%%%%%%%%%%%%%%%
    %%% Matemática:
    \input{static/utils/matematica}

    %%%%%%%%%%%%%%%%%%%%%%%%%%%%%%%%%%%%%%%%%%%%%%%%%%%%%%%%%%%%%%%%%%%%%%%%%%%%%%%%
    %%% Notas de rodapé:
    \input{static/utils/footnotes}

    %%%%%%%%%%%%%%%%%%%%%%%%%%%%%%%%%%%%%%%%%%%%%%%%%%%%%%%%%%%%%%%%%%%%%%%%%%%%%%%%
    %%% Referências cruzadas, links, citações:
    \input{static/utils/crossrefs}

    %%%%%%%%%%%%%%%%%%%%%%%%%%%%%%%%%%%%%%%%%%%%%%%%%%%%%%%%%%%%%%%%%%%%%%%%%%%%%%%%
    %%% Configurações para os hiperlinks em PDF:
    \input{static/utils/pdfconfig}

    %%%%%%%%%%%%%%%%%%%%%%%%%%%%%%%%%%%%%%%%%%%%%%%%%%%%%%%%%%%%%%%%%%%%%%%%%%%%%%%%
    %%% Glossário e índice remissivo:
    \input{static/utils/glossindex}

    %%%%%%%%%%%%%%%%%%%%%%%%%%%%%%%%%%%%%%%%%%%%%%%%%%%%%%%%%%%%%%%%%%%%%%%%%%%%%%%%
    %%% Referências bibliográficas:
    \input{static/utils/bibliografia}

    %%%%%%%%%%%%%%%%%%%%%%%%%%%%%%%%%%%%%%%%%%%%%%%%%%%%%%%%%%%%%%%%%%%%%%%%%%%%%%%%
    %%% Cores:
    \input{static/utils/cores}

    %%%%%%%%%%%%%%%%%%%%%%%%%%%%%%%%%%%%%%%%%%%%%%%%%%%%%%%%%%%%%%%%%%%%%%%%%%%%%%%%
    %%% Computação:
    \input{static/utils/computacao}

    %%%%%%%%%%%%%%%%%%%%%%%%%%%%%%%%%%%%%%%%%%%%%%%%%%%%%%%%%%%%%%%%%%%%%%%%%%%%%%%%
    %%% Gráficos:
    \input{static/utils/graficos}

    %%%%%%%%%%%%%%%%%%%%%%%%%%%%%%%%%%%%%%%%%%%%%%%%%%%%%%%%%%%%%%%%%%%%%%%%%%%%%%%%
    %%% Ícones:
    \input{static/utils/icones}

    %%%%%%%%%%%%%%%%%%%%%%%%%%%%%%%%%%%%%%%%%%%%%%%%%%%%%%%%%%%%%%%%%%%%%%%%%%%%%%%%
    %%% Blocos:
    \input{static/utils/blocos}

    %%%%%%%%%%%%%%%%%%%%%%%%%%%%%%%%%%%%%%%%%%%%%%%%%%%%%%%%%%%%%%%%%%%%%%%%%%%%%%%%
    %%% Boxes coloridos:
    \input{static/utils/colorbox}

    %%%%%%%%%%%%%%%%%%%%%%%%%%%%%%%%%%%%%%%%%%%%%%%%%%%%%%%%%%%%%%%%%%%%%%%%%%%%%%%%
    %%% Tabelas:
    \input{static/utils/tabelas}

    %%%%%%%%%%%%%%%%%%%%%%%%%%%%%%%%%%%%%%%%%%%%%%%%%%%%%%%%%%%%%%%%%%%%%%%%%%%%%%%%
    %%% Ambientes de listas:
    \input{static/utils/listas}

    %%%%%%%%%%%%%%%%%%%%%%%%%%%%%%%%%%%%%%%%%%%%%%%%%%%%%%%%%%%%%%%%%%%%%%%%%%%%%%%%
    %%% Ambientes floats e similares:
    \input{static/utils/floats}

    %%%%%%%%%%%%%%%%%%%%%%%%%%%%%%%%%%%%%%%%%%%%%%%%%%%%%%%%%%%%%%%%%%%%%%%%%%%%%%%%
    %%% Ajustes para a classe EXAM:
    \input{static/utils/exam}

    %%%%%%%%%%%%%%%%%%%%%%%%%%%%%%%%%%%%%%%%%%%%%%%%%%%%%%%%%%%%%%%%%%%%%%%%%%%%%%%%
    %%% Ajustes para textos:
    \input{static/utils/texto}

    %%%%%%%%%%%%%%%%%%%%%%%%%%%%%%%%%%%%%%%%%%%%%%%%%%%%%%%%%%%%%%%%%%%%%%%%%%%%%%%%
    %%% Ambientes, aliases, comandos e customizações em geral para serem
    %%% utilizadas nos documentos. Defina aqui qualquer customização específica
    %%% para seus documentos!
    \input{static/utils/customizacoes}

    %%%%%%%%%%%%%%%%%%%%%%%%%%%%%%%%%%%%%%%%%%%%%%%%%%%%%%%%%%%%%%%%%%%%%%%%%%%%%%%%
    %%% Ajuste do layout, espaçamento de linhas e etc.:
    \geometry{a4paper, portrait, left=2.0cm, right=2.0cm, top=2.0cm, bottom=2.0cm}
    \singlespacing

    %%%%%%%%%%%%%%%%%%%%%%%%%%%%%%%%%%%%%%%%%%%%%%%%%%%%%%%%%%%%%%%%%%%%%%%%%%%%%%%%
    %%% Configurações de cabeçalho e rodapé (não use fancyhdr pois ocorre conflito):
    \pagestyle{headandfoot}
    \runningheadrule
    \firstpageheader{}{}{}
    \runningheader{Sem disciplina}{}{Avaliação}
    \firstpagefooter{}{}{Página \thepage\ de \numpages}
    \runningfooter{}{}{Página \thepage\ de \numpages}
    \runningfootrule

    %%%%%%%%%%%%%%%%%%%%%%%%%%%%%%%%%%%%%%%%%%%%%%%%%%%%%%%%%%%%%%%%%%%%%%%%%%%%%%%%
    %%% Configurações para as propriedades do PDF:
    \hypersetup{
    hidelinks,           % Comente para web, descomente para impressão
    colorlinks=true,      % True para web, False para impressão
    pdftitle=PF: Avaliação,
    pdfauthor=Matheus Endlich Silveira,
    pdfsubject=Sem disciplina,
    pdfkeywords=['Sem tag'],
    pdfinfo={
        CreationDate={}, % Ex.: D:AAAAMMDDHH24MISS
        ModDate={}       % Ex.: D:AAAAMMDDHH24MISS
    }
    }
    \input{static/utils/pdfconfig}

    %%%%%%%%%%%%%%%%%%%%%%%%%%%%%%%%%%%%%%%%%%%%%%%%%%%%%%%%%%%%%%%%%%%%%%%%%%%%%%%%
    %%% Começa o documento
    \begin{document}

    %%%%%%%%%%%%%%%%%%%%%%%%%%%%%%%%%%%%%%%%%%%%%%%%%%%%%%%%%%%%%%%%%%%%%%%%%%%%%%%%
    %%% Folha de instruções padronizadas da UVV para a prova
    \pdfbookmark[1]{Instruções}{instrucoes}
    \begin{figure}[H]
        \begin{center}
            \includegraphics[scale=0.3]{{static/imgs/logo-uvv.png}}
        \end{center}	
    \end{figure}

    \vspace{-1cm}

    \begin{table}[H]
        \centering
        \begin{tabular}{|llll|}
            \hline
            \multicolumn{2}{|l|}{\textbf{Disciplina:} Sem disciplina} & 
            \multicolumn{1}{l|}{\multirow{3}{*}{\textbf{\makecell{Nota:\\ \,\\ \,}}}} & 
            \multirow{3}{*}{\textbf{\makecell{Coordenador:\\ \,\\ \,}}} \\ 
            \cline{1-2}
            \multicolumn{2}{|l|}{\textbf{Professor:} Matheus Endlich Silveira \hspace{4.72cm} } & 
            \multicolumn{1}{l|}{} & \\ 
            \cline{1-2}
            \multicolumn{2}{|l|}{\textbf{Aluno:}} & 
            \multicolumn{1}{l|}{} & \\ 
            \hline
            \multicolumn{1}{|l|}{\textbf{Turma:} \hspace{6.3cm} } & 
            \multicolumn{1}{l|}{\textbf{Semestre:}} & 
            \multicolumn{2}{l|}{\textbf{Valor:} 10 pontos} \\ 
            \hline
            \multicolumn{1}{|l|}{\textbf{Data:}} & 
            \multicolumn{3}{l|}{\textbf{Avaliação:}} \\ 
            \hline
        \end{tabular}
    \end{table}

    \begin{center}
        \fbox{\setlength\fboxsep{0.3cm} \fbox{\parbox{15.4cm}{
            \begin{center}
                \textbf{INSTRUÇÕES PARA A REALIZAÇÃO DESTA AVALIAÇÃO:}
            \end{center}
            \begin{itemize}[leftmargin=*]
                \item Esta é a primeira avaliação da disciplina Sem disciplina, 
                    e engloba o conteúdo referente Sem tag.
                \item Esta avaliação contém 30 questões objetivas, \textbf{todas obrigatórias}.
                \item \textbf{Leia atentamente cada questão!} As questões objetivas terão uma,
                    e apenas uma única, resposta a ser marcada.
                \item \textbf{Desligue o celular} antes de começar. A avaliação é
                    \textbf{individual} e \textbf{sem consulta}.
                \item As questões podem ser respondidas, na prova, com lápis ou caneta. Ao final
                    da prova \textbf{{VOCÊ DEVERÁ PREENCHER O GABARITO COM CANETA
                    DE COR PRETA}} \textbf{\underline{OU AZUL}} para que seja feita a correção
                    automatizada através da leitura do padrão de respostas marcado no
                    gabarito.
                \item \textbf{CUIDADO ao preencher o gabarito} pois eles são \textbf{INDIVIDUAIS
                    e NOMEADOS} (cada estudante tem um gabarito próprio específico, com um
                    código de identificação único). \textbf{Questões rasuradas são anuladas}
                    pelo software de leitura do gabarito.
                \item Ao final da prova, \textbf{DEVOLVA A PROVA E O GABARITO} para o professor.
                \item Siga todas as normas de \textbf{Integridade Acadêmica} da disciplina.
                    Alunos que forem flagrados com qualquer espécie de ``cola'' ou trocando
                    informações com outros alunos terão suas avaliações recolhidas, as notas
                    zeradas, e a situação será encaminhada para a coordenação para a aplicação
                    das penalidades previstas pela universidade.
                \item Com 90 minutos de avaliação você terá tempo de até 3,min por questão,
                    em média.
                \item Boa avaliação!
            \end{itemize}
            \vspace{2cm}
        }}}
    \end{center}
    \newpage

    %%%%%%%%%%%%%%%%%%%%%%%%%%%%%%%%%%%%%%%%%%%%%%%%%%%%%%%%%%%%%%%%%%%%%%%%%%%%%%%%
    %%% Inicia o bloco de questões:
    \begin{questions}

    \newpage
    \question

Dentre as definições a seguir, ligadas ao conceito de visões do modelo relacional, qual delas é \textbf{INCORRETA}?


\begin{choices}
\choice Uma visão relacional é uma relação virtual, derivada de relações base a partir da especificação de operações da álgebra relacional.
\choice Programas aplicativos do banco de dados podem ser executados sobre visões de relações da base de dados.
\choice Uma visão é útil por representar uma percepção particular do banco de dados, compartilhado por muitos aplicativos.
\choice O gerenciamento de visões envolve a conversão da consulta do usuário sobre as visões para a consulta sobre as relações base.
\choice Uma visão relacional é uma relação virtual que nunca é materializada.
\end{choices}

\question

São modelos de dados legados, pré-relacionais, exceto:


\begin{choices}
\choice Hierárquico
\choice Objeto
\choice Nenhuma das alternativas anteriores
\choice Rede
\choice Flat file
\end{choices}

\question

Um sistema de gerenciamento de bancos de dados (SGBD) é um software servidor

que nos permite definir, construir, manipular, integrar e compartilhar bancos

de dados entre diversos usuários e aplicações. São exemplos de

SGBD, \textbf{exceto}:
\begin{choices}
\choice PostgreSQL
\choice Oracle SQL Developer
\choice SQL Server
\choice MongoDB
\choice Oracle Database
\end{choices}

\question

A derivada da função \( f(x) = x^x \) é igual a:
\begin{choices}
\choice \( x^x \)
\choice \( x^x \ln(x) \)
\choice \( x^x (\ln(x) + 1) \)
\choice \( x^x (\ln(x) + x) \)
\choice \( x x^{x-1} \)
\end{choices}

\question

Histograma de uma imagem com \( K \) tons de cinza é:
\begin{choices}
\choice Contagem dos pixels da imagem.
\choice Nenhuma alternativa acima.
\choice Contagem do número de vezes que cada um dos \( K \) tons de cinza ocorreu na imagem.
\choice Contagem do número de objetos encontrados na imagem.
\choice Contagem do número de tons de cinza que ocorreram na imagem.
\end{choices}

\question

Qual o significado de coerência de memórias \textit{cache} em sistemas multiprocessados?


\begin{choices}
\choice Caches em processadores diferentes sempre lêem os mesmos dados ao mesmo tempo.
\choice Caches em processadores diferentes nunca interagem entre si.
\choice Caches em processadores diferentes sempre contêm o mesmo dado válido para a mesma linha de cache.
\choice Caches em processadores diferentes podem possuir dados diferentes associados à mesma linha de cache.
\choice Caches em processadores diferentes nunca compartilham a mesma linha de cache.
\end{choices}

\question

Considere as seguintes afirmativas:

\begin{itemize}

    \item O modelo matemático de uma lista é a sequência linear de itens, cuja principal propriedade estrutural é a posição relativa dos elementos dentro da sequência.

    \item A fila e a pilha são consideradas casos especiais da lista.

    \item Numa fila a inserção e a retirada são feitas no mesmo extremo.

    \item Numa lista a inserção e a retirada podem ser feitas em qualquer posição.

    \item Numa pilha apenas a inserção pode ser feita em qualquer posição.

\end{itemize}



Quais são as afirmativas verdadeiras?
\begin{choices}
\choice somente II, III e IV
\choice somente I, II e IV
\choice somente I e III
\choice todas
\choice somente II, IV e V
\end{choices}

\question

Supondo a Relação $\texttt{PROJ (PNO, Orçam)}$, com chave primária $\texttt{PNO}$, a Relação $\texttt{EMP (ENO, ENome, Cargo)}$ com chave primária $\texttt{ENO}$, e a Relação $\texttt{DSG (ENO, PNO, Dur, Resp)}$, com chave primária \{\texttt{ENO, PNO}\}, chave estrangeira $\texttt{PNO}$ em relação a $\texttt{PROJ}$ e chave estrangeira $\texttt{ENO}$ em relação a $\texttt{EMP}$. Qual das expressões da álgebra relacional abaixo $\textbf{NÃO}$ corresponde à seguinte consulta SQL:



\begin{verbatim}

SELECT ENome

FROM EMP, PROJ, DSG

WHERE EMP.ENO = DSG.ENO

      AND PROJ.PNO = DSG.PNO

      AND Dur > 36

\end{verbatim}
\begin{choices}
\choice \( \pi_{ENome} (\sigma_{Dur > 36} ((\pi_{ENome, ENO} (\text{EMP})) \Join_{ENO} (\text{PROJ} \Join_{PNO} (\text{DSG})))) \)
\choice \( \pi_{ENome} (\text{PROJ} \Join_{PNO} (\text{EMP} \Join_{ENO} \sigma_{Dur > 36} (\text{DSG}))) \)
\choice \( \pi_{ENome} (\sigma_{Dur > 36} ((\pi_{PNO} (\text{PROJ})) \Join_{PNO} (\text{EMP} \Join_{ENO} (\pi_{Dur} (\text{DSG})))) \)
\choice \( \pi_{ENome} (\text{PROJ} \Join_{PNO} (\text{EMP} \Join_{ENO} \sigma_{Dur > 36} (\pi_{Dur} (\text{DSG})))) \)
\choice \( \pi_{ENome} (\text{PROJ} \Join_{PNO} (\sigma_{Dur > 36} (\text{EMP} \Join_{ENO} (\text{DSG})))) \)
\end{choices}

\question

Sobre a UML, quais das seguintes afirmações são verdadeiras?



\begin{itemize}

    \item[I)] A UML é o método de desenvolvimento de software mais utilizado na atualidade.

    \item[II)] A UML é uma evolução das linguagens para especificação dos conceitos dos métodos de Booch, OMT e OOSE e também de outros métodos de especificação de requisitos de software orientados a objetos ou não.

    \item[III)] A UML é composta dos seguintes diagramas: Diagrama de Caso de Uso, Diagrama de Classes, Diagrama de Colaboração, Diagrama de Estados, entre outros.

    \item[IV)] Em UML pode-se representar tão somente relacionamentos de Agregação, Associação e Composição.

\end{itemize}
\begin{choices}
\choice Apenas as alternativas II e III.
\choice Apenas as alternativas I, II e III.
\choice Todas as alternativas.
\choice Nenhuma delas.
\choice Apenas as alternativas III e IV.
\end{choices}

\question

Um compilador detecta:
\begin{choices}
\choice todos os erros citados acima
\choice erros aritméticos
\choice erros de sintaxe do programa 
\choice erros nos resultados gerados pelo programa
\choice erros que podem ocorrer durante a execução do programa
\end{choices}

\question

Quais dos algoritmos de ordenação abaixo possuem tempo no pior caso e tempo médio de execução proporcional a \( O(n \log n) \).
\begin{choices}
\choice Bubble sort e Quick sort
\choice Quicksort e merge sort
\choice Merge sort e bubble sort
\choice Merge sort e heap sort
\choice Heap sort e selection sort
\end{choices}

\question

A arquitetura ilustrada na figura abaixo é um exemplo típico de qual sistema?

(CU = unidade de controle; IS = \ingles{stream} de instruções; PU = unidade de

processamento; DS = \ingles{stream} de dados; LM = memória local)

\begin{figure}[H]

\begin{center}

\includegraphics[width=0.5\linewidth]{static/imgs/defaul.png}

\end{center}

\end{figure}

\vspace{-1cm}
\begin{choices}
\choice Processadores multicore
\choice Clusters
\choice Multiprocessadores simétricos
\choice Acesso não uniforme à memória
\choice Processadores seqüenciais
\end{choices}

\question

Uma definição lógica, abstrata e auto-contida dos objetos que modelam a

estrutura de armazenamento de dados, dos operadores que modelam o comportamento

e dos demais conceitos que modelam a integridade, e que, juntos, constituem a

máquina abstrata com a qual os usuários interagem é chamada de:


\begin{choices}
\choice Nenhuma das respostas acima
\choice Modelo entidade-relacionamento
\choice Modelo relacional
\choice Modelo de dados
\choice Banco de dados
\end{choices}

\question

Seu chefe que fazer uma campanha de venda para produtos que custam entre 20 e 30

reais, mas está preocupado porque ouviu de um dos vendedores que diversos

produtos nessa faixa de preço estão com o estoque zerado. Para verificar a

situação seu chefe pediu para que você prepare um relatório que mostre os

produtos que custam entre 20 e 30 reais e que estão com o estoque zerado. Você

preparou o seguinte relatório para seu chefe:

\begin{tcolorbox}

 \begin{verbatim}1  SELECT l.nome AS loja

2       , p.nome AS produto

3       , p.preco_unitario AS preco

4       , e.quantidade AS qtd

5  FROM produtos p

6  INNER JOIN estoques e ON (e.produto_id = p.produto_id)

7  INNER JOIN lojas l ON (e.loja_id = l.loja_id)

8  WHERE (p.preco_unitario >= 20 OR p.preco_unitario <= 30)

9    AND e.quantidade = 0

10 ORDER BY loja, produto;

 \end{verbatim}

\end{tcolorbox}

O seu relatório:
\begin{choices}
\choice Está errado pois a tabela que deveria estar na cláusula FROM é a tabela ``lojas', não a tabela ``produtos'.
\choice Está correto e trará exatamente, os produtos com o estoque zerado e na faixa de preço requisitada pelo seu chefe. Além disso ordena o resultado pelo nome da loja e pelo nome do produto.
\choice Está errado pois vai trazer também produtos fora da faixa de preço que o seu chefe quer.
\choice Nenhuma das alternativas anteriores.
\choice Está errado pois há um erro de sintaxe nas linhas 6 e 7.
\end{choices}

\question

Na linguagem SQL, qual das seguintes palavras-chave é usada para modificar os

dados em uma tabela existente?


\begin{choices}
\choice INSERT
\choice CREATE
\choice Nenhuma das alternativas acima
\choice UPDATE
\choice SELECT
\end{choices}

\question

O determinante da matriz dada abaixo é

$\begin{pmatrix}

2 & 7 & 9 & -1 & 1 \\

2 & 8 & 3 & 1 & 0 \\

-1 & 0 & 4 & 3 & 0 \\

2 & 0 & -1 & 0 & 0 \\

3 & 0 & 0 & 0 & 0 \\

\end{pmatrix}$
\begin{choices}
\choice 86
\choice -86
\choice 96
\choice 46
\choice -96
\end{choices}

\question

A medida da interconexão entre os módulos de uma estrutura de software é denominada e que também é usada em projetos orientados a objetos é: 
\begin{choices}
\choice acoplamento
\choice unidade funcional
\choice coesão
\choice abstração procedimental
\choice ocultamento da informação
\end{choices}

\question

Filtro da mediana é:
\begin{choices}
\choice Indicado para detectar tonalidades específicas em uma imagem.
\choice Indicado para detectar formas específicas em imagens.
\choice Indicado para atenuar ruído com preservação de bordas (i.e., rápidas transições de nível em imagens).
\choice Nenhuma das respostas acima.
\choice Indicado para detectar bordas em imagens.
\end{choices}

\question

Qual das seguintes afirmações sobre crescimento assintótico de funções não é verdadeira:
\begin{choices}
\choice Se \( f(n) = O(g(n)) \) então \( g(n) = O(f(n)) \)
\choice Se \( f(n) = O(g(n)) \) e \( g(n) = O(h(n)) \) então \( f(n) = O(h(n)) \)
\choice \( \log n^2 = O(\log n) \)
\choice \( 2n^2 + 3n + 1 = O(n^2) \)
\choice \( 2^{n+1} = O(2^n) \)
\end{choices}

\question

São motivos importantes e gerais para o estudo das linguagens de programação,

exceto:
\begin{choices}
\choice Avanço geral da computação
\choice Melhorar o embasamento na escolha das linguagens para um projeto
\choice Entender a importância da implementação
\choice Aprender a sintaxe e a semântica de uma determinada linguagem
\choice Aumentar a capacidade de expressar idéias
\end{choices}

\question

Marque a alternativa onde todos os conceitos estão corretos.


\begin{choices}
\choice Nenhuma das alternativas anteriores.
\choice Em um diagrama de fluxo de dados uma entidade externa representa uma fonte ou destino das informações processadas pelo sistema e está fora dos limites do sistema modelado; cada processo pode ser refinado, para explicitar um maior detalhamento; um DFD pode conter vários níveis de detalhamento; um processo é um transformador de informação; o nível 0 de um DFD representa o sistema como um todo e indica as principais fontes e destinos das informações, usualmente referenciado por Diagrama de Contexto.
\choice Em um diagrama de fluxo de dados uma entidade externa representa um produtor ou um consumidor de informação e está fora dos limites do sistema modelado; cada processo deve ser refinado, para explicitar um maior detalhamento; um DFD pode conter vários níveis de detalhamento; um processo é um transformador de informação e também está fora do sistema; o nível 0 de um DFD representa o sistema como um todo e indica os principais usuários e as funções do sistema.
\choice Em um diagrama de fluxo de dados, uma entidade externa representa um produtor ou um consumidor de informação e está fora dos limites do sistema modelado; cada processo pode ser refinado, para explicitar um maior detalhamento; um DFD contém dois níveis de detalhamento; um processo é um transformador de informação e também está fora do sistema; o nível 0 de um DFD representa o sistema como um todo e indica os principais usuários e as funções do sistema.
\choice Em um diagrama de fluxo de dados uma entidade externa representa uma fonte ou destino das informações processadas pelo sistema e está fora dos limites do sistema modelado; cada processo pode ser refinado, para explicitar um maior detalhamento; um DFD pode conter vários níveis de detalhamento; um processo é um transformador de informação e também está fora do sistema; o nível 0 de um DFD representa o sistema como um todo e indica as principais fontes e destinos das informações.
\end{choices}

\question

A média aritmética de uma lista de 50 números é 50. Se dois desses números, 51 e 97, forem suprimidos dessa lista a média dos restantes será


\begin{choices}
\choice 40
\choice 51
\choice 49
\choice 47
\choice 50
\end{choices}

\question

São vantagens da utilização de threads no espaço do usuário, exceto:
\begin{choices}
\choice O escalonamento pode ser específico para a aplicação.
\choice O sistema operacional escalona a thread.
\choice A criação e o gerenciamento das threads são mais eficientes.
\choice Nenhuma modificação é necessária no kernel.
\choice Maior portabilidade da aplicação.
\end{choices}

\question

Dentre as definições a seguir, ligadas ao conceito de normalização do modelo relacional, qual delas é \textbf{INCORRETA}?


\begin{choices}
\choice As formas normais se baseiam em certas estruturas de dependências.
\choice A normalização é um processo passo a passo reversível de substituição de uma dada coleção de relações por sucessivas coleções de relações as quais possuem uma estrutura progressivamente mais simples e mais regular.
\choice A primeira forma normal estabelece que os atributos da relação contêm apenas valores atômicos.
\choice As relações que obedecem à primeira forma normal não apresentam anomalias.
\choice O objetivo da normalização é eliminar várias anomalias (ou aspectos indesejáveis) de uma relação.
\end{choices}

\question

No que diz respeito as vantagens da arquitetura de micro-núcleo para sistemas operacionais em relação a arquiteturas de núcleo monolítico, quais das seguintes afirmações são verdadeiras?



\begin{itemize}

    \item[I] A arquitetura de micro-núcleo facilita a depuração do SO.

    \item[II] A arquitetura de micro-núcleo permite um número menor de mudanças de contexto.

    \item[III] A arquitetura de micro-núcleo facilita a reconfiguração de serviços do SO pois a maioria deles reside em espaço de usuário.

\end{itemize}
\begin{choices}
\choice Todas são verdadeiras
\choice II e III
\choice I e II
\choice I e III
\choice Apenas I
\end{choices}

\question

Considere o algoritmo da busca sequencial de um elemento em um conjunto com \( n \) elementos. A expressão que representa o tempo médio de execução desse algoritmo para uma busca bem sucedida é:
\begin{choices}
\choice \( n^2 \)
\choice \( (n + 1)/2 \)
\choice \( \log_2 n \)
\choice \( n \log n \)
\choice \( n(n + 1)/2 \)
\end{choices}

\question

Considerando a rede de Petri abaixo, quais das alternativas são verdadeiras?



\begin{itemize}

    \item[I)] O lugar \( A \) está habilitado a disparar.

    \item[II)] Apenas a transição \( T1 \) está habilitada a disparar.

    \item[III)] A sequência de transições \( (T1, T2, T3, T2) \) pode ser disparada, nessa ordem.

    \item[IV)] A transição \( T4 \) nunca poderá ser disparada.

\end{itemize}



\begin{figure}[H]

    \centering

    \includegraphics[width=0.5\linewidth]{static/imgs/defaul.png}

    \caption{Questão 49}

    \label{fig:enter-label}

\end{figure}
\begin{choices}
\choice Todas as alternativas.
\choice Apenas as alternativas III, IV.
\choice Apenas as alternativas II e III.
\choice Apenas as alternativas I e III.
\choice Apenas as alternativas II, III e IV.
\end{choices}

\question

Assinale quantas sequências de caracteres a seguir são reconhecidas pelo autômato finito abaixo. As quatro sequências de caracteres (separadas por vírgulas) são: \(0\), \(+567\), \(-89.5\), \(-3e3\).



\begin{figure} [H]

    \centering

    \includegraphics[width=0.5\linewidth]{static/imgs/defaul.png}

    \caption{Questão 25}

    \label{fig:enter-label}

\end{figure}



Qual das seguintes alternativas indica o número de sequências reconhecidas pelo autômato?
\begin{choices}
\choice 2
\choice 3
\choice 1
\choice 4
\choice 0
\end{choices}

\question

(ENADE 2014) Considere o seguinte banco de dados ilustrado no diagrama relaconal

abaixo:

\begin{figure}[H]

  \centering

  \caption{Estados e fornecedores}

  \label{fig:estados}

  %\fbox{

    \includegraphics[width=0.5\linewidth]{static/imgs/defaul.png}

  %}\\

  %\footnotesize{Fonte: xxxx.}

\end{figure}

A expressão SQL que obtém os nomes dos estados para os quais não há fornecedores

cadastrados é:


\begin{choices}
\choice \begin{verbatim}SELECT e.nome

FROM estado e,

FROM fornecedor f

WHERE e.uf = f.uf;

\end{verbatim}


\choice \begin{verbatim}SELECT e.nome

FROM estado e

WHERE e.uf IN (SELECT f.uf

               FROM fornecedor f);

\end{verbatim}


\choice \begin{verbatim}SELECT e.nome

FROM estado e

WHERE e.uf NOT IN (SELECT f.uf

                   FROM fornecedor f);

\end{verbatim}


\choice \begin{verbatim}SELECT e.uf

FROM estado e

WHERE e.nome NOT IN (SELECT f.uf

                     FROM fornecedor f);

\end{verbatim}
\choice \begin{verbatim}SELECT e.nome

FROM estado e,

FROM fornecedor f

WHERE e.nome = f.uf;

\end{verbatim}


\end{choices}

\question

Quanto ao TCP, é INCORRETO afirmar:
\begin{choices}
\choice Usa janelas deslizantes para implementar o controle de fluxo e erro.
\choice É um protocolo orientado a conexão.
\choice Possui um campo de \textit{checksum} que valida as informações de seu cabeçalho, mas não valida as informações de \textit{payload} (campo de dados).
\choice É um protocolo do nível de transporte.
\choice Utiliza portas para permitir a comunicação entre processos localizados em dispositivos diferentes.
\end{choices}



    %%%%%%%%%%%%%%%%%%%%%%%%%%%%%%%%%%%%%%%%%%%%%%%%%%%%%%%%%%%%%%%%%%%%%%%%%%%%%%%%
    %%% Finaliza o bloco de questões:
    \end{questions}
    \end{document}
    